\documentclass[11pt]{article}

\usepackage[margin=1in]{geometry}
\usepackage{graphicx}
\usepackage{amsmath}
\usepackage{amssymb}
\usepackage{booktabs}
\usepackage{hyperref}
\usepackage{float}
\usepackage{subcaption}
\usepackage{multirow}
\usepackage{enumitem}

\title{
Residual Universality Manifolds:\\
Boundary Structure, Geodesic Transport, and Diffusion Persistence
}

\author{
Dan Hawkley \\
Independent Researcher \\
Boulder, Colorado \\
\texttt{github.com/thinkthoughts/residual-phase-lock}
}

\date{May 2026}

\begin{document}

\maketitle

\begin{abstract}

Residual embedding systems are often treated as high-dimensional statistical artifacts with limited geometric interpretation. In this work, we study residual universality manifolds constructed from multiple graph topology families and investigate whether persistent geometric structure survives across topology scales and transitions.

We analyze boundary confidence structure, nearest-family partitioning, geodesic transport paths, diffusion persistence, and spectral organization across residual manifolds generated from Erdős--Rényi, modular clustered, ring lattice, scale free, and small world systems.

The experiments show that residual manifolds exhibit structured density organization, nontrivial transport geometry, diffusion persistence regimes, and measurable topology separability. Geodesic transport exposes bridgeability relationships between topology families, while diffusion and spectral analyses reveal persistent manifold organization under perturbation and scale variation.

These results suggest that residual embeddings contain interpretable geometric structure that can be probed through transport, diffusion, and spectral methods rather than treated solely as unstructured residual spaces.

\end{abstract}

\section{Introduction}

Residual representations appear throughout machine learning, graph systems, dynamical systems, and embedding pipelines. Despite their widespread use, residual embeddings are often interpreted primarily through statistical performance metrics rather than through geometric or topological organization.

This work studies whether structured manifold geometry persists across residual systems generated from multiple topology families. Rather than treating residual embeddings as unordered point clouds, we investigate whether measurable geometric continuity survives across:

\begin{itemize}[leftmargin=1.5em]
    \item topology families,
    \item graph scales,
    \item transport paths,
    \item diffusion processes,
    \item and spectral decompositions.
\end{itemize}

The experiments focus on five graph topology families:

\begin{itemize}[leftmargin=1.5em]
    \item Erdős--Rényi,
    \item modular clustered,
    \item ring lattice,
    \item scale free,
    \item small world.
\end{itemize}

Residual embeddings are projected into low-dimensional manifolds where density structure, transport paths, diffusion persistence, and Laplacian spectral organization can be analyzed directly.

The central question is whether residual universality manifolds contain persistent geometric structure that survives topology variation and perturbation.

\section{Residual Manifold Construction}

Residual manifolds are constructed from graph topology families generated across multiple scales. Each graph family contributes a sequence of embeddings associated with increasing graph size and connectivity variation.

For each generated graph:

\begin{enumerate}[leftmargin=1.5em]
    \item graph statistics are computed,
    \item residual feature vectors are assembled,
    \item embeddings are normalized,
    \item dimensionality reduction is applied,
    \item manifold coordinates are extracted.
\end{enumerate}

The resulting residual manifold provides a geometric representation of topology families embedded into a shared latent space.

The manifold coordinates are then used for:

\begin{itemize}[leftmargin=1.5em]
    \item density estimation,
    \item nearest-family partitioning,
    \item geodesic transport analysis,
    \item diffusion simulation,
    \item spectral decomposition.
\end{itemize}

\section{Boundary Structure and OOD Regions}

We first investigate density structure and nearest-family organization across the residual manifold.

\subsection{Residual Density Structure}

Figure~\ref{fig:density_manifold} shows continuous density structure over the known residual manifold.

\begin{figure}[H]
\centering
\includegraphics[width=0.9\linewidth]{../figures/24_density_known_manifold.png}
\caption{
Continuous density structure over the residual universality manifold. Density regions reveal persistent organization rather than uniform embedding dispersion.
}
\label{fig:density_manifold}
\end{figure}

Distinct topology families occupy partially separable regions while maintaining smooth continuity transitions between neighboring structures.

\subsection{Nearest-Family Boundary Structure}

Figure~\ref{fig:decision_field} shows nearest-family decision regions across the manifold.

\begin{figure}[H]
\centering
\includegraphics[width=0.9\linewidth]{../figures/24_decision_field_nearest_family.png}
\caption{
Nearest-family residual manifold partitioning. Topology families remain partially separable while exhibiting continuity regions between neighboring structures.
}
\label{fig:decision_field}
\end{figure}

Boundary transitions reveal that topology families are neither fully collapsed nor fully isolated. Instead, residual manifolds contain structured continuity regions connecting related topology systems.

\subsection{Universality Potential Field}

Figure~\ref{fig:universality_potential} shows the universality potential field across the residual manifold.

\begin{figure}[H]
\centering
\includegraphics[width=0.9\linewidth]{../figures/24_universality_potential_field.png}
\caption{
Residual universality potential field. Smooth transitions indicate persistent geometric continuity across topology families.
}
\label{fig:universality_potential}
\end{figure}

The field structure suggests that residual manifolds contain coherent geometric organization beyond isolated clustering behavior.

\section{Geodesic Transport and Bridgeability}

We next study transport structure between topology families.

\subsection{Geodesic Transport Paths}

Figure~\ref{fig:geodesic_paths} shows geodesic transport paths across the manifold.

\begin{figure}[H]
\centering
\includegraphics[width=0.95\linewidth]{../figures/25_geodesic_transport_paths.png}
\caption{
Geodesic transport paths across topology families. Paths reveal structured bridgeability relationships within the residual manifold.
}
\label{fig:geodesic_paths}
\end{figure}

Some topology families exhibit low-cost transitions, while others remain separated by higher transport curvature and bottleneck structure.

\subsection{Bridgeability Matrix}

Figure~\ref{fig:bridge_matrix} shows transport costs and bridgeability relationships.

\begin{figure}[H]
\centering
\includegraphics[width=0.95\linewidth]{../figures/25_bridge_matrix_heatmap.png}
\caption{
Geodesic transport cost and bridgeability relationships between topology families.
}
\label{fig:bridge_matrix}
\end{figure}

The matrix reveals asymmetries in manifold connectivity and identifies topology families that act as transport bridges between otherwise separated regions.

\section{Diffusion Persistence and Spectral Organization}

We next investigate diffusion stability and spectral persistence across the manifold.

\subsection{Residual Diffusion Graph}

Figure~\ref{fig:diffusion_graph} shows the residual diffusion kNN graph.

\begin{figure}[H]
\centering
\includegraphics[width=0.95\linewidth]{../figures/26_diffusion_knn_graph.png}
\caption{
Residual diffusion kNN graph across topology families and graph scales.
}
\label{fig:diffusion_graph}
\end{figure}

Diffusion connectivity preserves partial topology separability while allowing controlled transport between neighboring structures.

\subsection{Entropy Growth and Persistence Decay}

Figure~\ref{fig:entropy_persistence} shows entropy growth and persistence decay under diffusion.

\begin{figure}[H]
\centering
\includegraphics[width=0.95\linewidth]{../figures/26_entropy_and_persistence.png}
\caption{
Diffusion entropy growth and universality persistence decay across topology families.
}
\label{fig:entropy_persistence}
\end{figure}

Different topology families exhibit distinct persistence regimes. Some systems diffuse rapidly into shared manifold regions, while others preserve localized structure over longer diffusion times.

\subsection{Phase Transition Structure}

Figure~\ref{fig:phase_transition} shows diffusion phase transition structure.

\begin{figure}[H]
\centering
\includegraphics[width=0.9\linewidth]{../figures/26_phase_transition_diagram.png}
\caption{
Universality persistence phase transitions under diffusion parameter variation.
}
\label{fig:phase_transition}
\end{figure}

The transition regions reveal measurable collapse thresholds where manifold persistence becomes unstable.

\subsection{Spectral Modes}

Figure~\ref{fig:spectral_modes} shows Laplacian spectral mode organization.

\begin{figure}[H]
\centering
\includegraphics[width=0.95\linewidth]{../figures/26_spectral_modes.png}
\caption{
Residual manifold Laplacian spectral modes.
}
\label{fig:spectral_modes}
\end{figure}

Spectral organization further supports the existence of persistent manifold structure rather than random embedding dispersion.

\section{Discussion}

The experiments demonstrate that residual manifolds exhibit measurable geometric structure across multiple independent probes:

\begin{itemize}[leftmargin=1.5em]
    \item density continuity,
    \item nearest-family partitioning,
    \item geodesic transport,
    \item diffusion persistence,
    \item spectral organization.
\end{itemize}

The residual manifolds are neither fully collapsed nor completely disconnected. Instead, topology families remain partially separable while participating in structured continuity relationships.

Transport geometry reveals bridgeability between topology families, while diffusion persistence exposes stability regimes and collapse thresholds across the manifold.

These results suggest that residual embeddings contain interpretable geometric organization that can be analyzed through manifold methods rather than treated solely as statistical residual artifacts.

\section{Conclusion}

This work studied residual universality manifolds across multiple topology families using transport, diffusion, and spectral probes.

The experiments showed that residual embeddings contain persistent geometric structure characterized by:

\begin{itemize}[leftmargin=1.5em]
    \item structured density organization,
    \item continuity boundaries,
    \item geodesic bridgeability,
    \item diffusion persistence,
    \item spectral phase structure.
\end{itemize}

Residual manifolds therefore provide a measurable geometric framework for analyzing continuity and organization across structured embedding systems.

\bibliographystyle{plain}

\begin{thebibliography}{9}

\bibitem{diffusion}
R. R. Coifman and S. Lafon,
\textit{Diffusion Maps},
Applied and Computational Harmonic Analysis,
2006.

\bibitem{laplacian}
Fan Chung,
\textit{Spectral Graph Theory},
American Mathematical Society,
1997.

\bibitem{manifold}
Mikhail Belkin and Partha Niyogi,
\textit{Laplacian Eigenmaps for Dimensionality Reduction and Data Representation},
Neural Computation,
2003.

\bibitem{topology}
Albert-László Barabási,
\textit{Network Science},
Cambridge University Press,
2016.

\end{thebibliography}

\end{document}
